% GENERATED FILE. Edit canonical lessons, then run npm run generate:books.
% canonical-source-hash: fnv1a64:f21aab5b74a3306f
% canonical-lessons: LA-C231-verto, LA-C231-fruor, LA-C231-fingo, LA-C231-queror, LA-C231-desum

\chapter{To turn, To enjoy, To shape, To complain, To be missing}
\label{ch:la-to-turn-to-enjoy-to-shape-to-complain-to-be-missing}
\hlchaptermodality{231}

\begin{chapteropening}
\textbf{By the end of this chapter:} \emph{I can say to turn, to enjoy, to shape, to complain and to be missing.}

Complete the last lesson of chapter 231: I can say to turn, to enjoy, to shape, to complain and to be missing.
\end{chapteropening}

\section[vertō, vertere]{vertō, vertere — to turn}

\label{lesson:LA-C231-verto}



\begin{quote}
Before the new one: say the Latin for to announce, then the Latin for to appear.
\end{quote}

\subsection*{You'll want to know: vertō, vertere}
\textbf{vertō, vertere} — \textquotedblleft{}to turn\textquotedblright{}.

\begin{cousinweb}
\textbf{vertō} is inherited, usually traced to Proto-Indo-European \emph{*wert-}, \textquotedblleft{}to turn\textquotedblright{} (EtymDB 2.1); Sanskrit \emph{vart-}, \textquotedblleft{}to turn\textquotedblright{}, is a cousin (Lewis \& Short), and so is English \emph{-ward} in words like toward.
\end{cousinweb}

\begin{grammarlens}[title={What you've built}]
One more word for this chapter.
\end{grammarlens}

\subsection*{Your turn}
\begin{itemize}
\raggedright
  \item \emph{Say it:} \emph{vertō, vertere}
  \item \emph{Say it:} \emph{vertō, vertere}, once more
  \item \emph{Say it:} say \emph{vertō, vertere}
  \item \emph{Recall:} say the Latin for to announce, then the Latin for to appear, then say \emph{vertō, vertere} again
\end{itemize}

\begin{tcolorbox}[breakable,colback=teal!4,colframe=teal!35!black,title={Before you move on}]
What does vertō, vertere mean? (\textquotedblleft{}To turn\textquotedblright{}.) Say it once more.
\end{tcolorbox}
\section[fruor, fruī]{fruor, fruī — to enjoy}

\label{lesson:LA-C231-fruor}



\begin{quote}
Before the new one: say the Latin for to appear, then the Latin for to turn.
\end{quote}

\subsection*{You'll want to know: fruor, fruī}
\textbf{fruor, fruī} — \textquotedblleft{}to enjoy\textquotedblright{}.

\begin{cousinweb}
\textbf{fruor} is inherited, usually traced to Proto-Indo-European \emph{*bhrug-}, \textquotedblleft{}to make use of, to enjoy\textquotedblright{} (EtymDB 2.1). English \emph{fruit} comes from its family.
\end{cousinweb}

\begin{grammarlens}[title={What you've built}]
One more word for this chapter.
\end{grammarlens}

\subsection*{Your turn}
\begin{itemize}
\raggedright
  \item \emph{Say it:} \emph{fruor, fruī}
  \item \emph{Say it:} \emph{fruor, fruī}, once more
  \item \emph{Say it:} say \emph{fruor, fruī}
  \item \emph{Recall:} say the Latin for to appear, then the Latin for to turn, then say \emph{fruor, fruī} again
\end{itemize}

\begin{tcolorbox}[breakable,colback=teal!4,colframe=teal!35!black,title={Before you move on}]
What does fruor, fruī mean? (\textquotedblleft{}To enjoy\textquotedblright{}.) Say it once more.
\end{tcolorbox}
\section[fingō, fingere]{fingō, fingere — to shape, to imagine}

\label{lesson:LA-C231-fingo}



\begin{quote}
Before the new one: say the Latin for to turn, then the Latin for to enjoy.
\end{quote}

\subsection*{You'll want to know: fingō, fingere}
\textbf{fingō, fingere} — \textquotedblleft{}to shape, to imagine\textquotedblright{}.

\begin{cousinweb}
\textbf{fingō} is inherited, usually traced to Proto-Indo-European \emph{*dheygh-}, \textquotedblleft{}to mould, to shape\textquotedblright{} (EtymDB 2.1). English \emph{dough} is a cousin, and English \emph{fiction} and \emph{figure} come from its family.
\end{cousinweb}

\begin{grammarlens}[title={What you've built}]
One more word for this chapter.
\end{grammarlens}

\subsection*{Your turn}
\begin{itemize}
\raggedright
  \item \emph{Say it:} \emph{fingō, fingere}
  \item \emph{Say it:} \emph{fingō, fingere}, once more
  \item \emph{Say it:} say \emph{fingō, fingere}
  \item \emph{Recall:} say the Latin for to turn, then the Latin for to enjoy, then say \emph{fingō, fingere} again
\end{itemize}

\begin{tcolorbox}[breakable,colback=teal!4,colframe=teal!35!black,title={Before you move on}]
What does fingō, fingere mean? (\textquotedblleft{}To shape, to imagine\textquotedblright{}.) Say it once more.
\end{tcolorbox}
\section[queror, querī]{queror, querī — to complain}

\label{lesson:LA-C231-queror}



\begin{quote}
Before the new one: say the Latin for to enjoy, then the Latin for to shape.
\end{quote}

\subsection*{You'll want to know: queror, querī}
\textbf{queror, querī} — \textquotedblleft{}to complain\textquotedblright{}.

\begin{cousinweb}
\textbf{queror} is inherited, usually traced to Proto-Indo-European \emph{*kwes-}, \textquotedblleft{}to puff, to sigh\textquotedblright{} (EtymDB 2.1). English \emph{wheeze} is a cousin.
\end{cousinweb}

\begin{grammarlens}[title={What you've built}]
One more word for this chapter.
\end{grammarlens}

\subsection*{Your turn}
\begin{itemize}
\raggedright
  \item \emph{Say it:} \emph{queror, querī}
  \item \emph{Say it:} \emph{queror, querī}, once more
  \item \emph{Say it:} say \emph{queror, querī}
  \item \emph{Recall:} say the Latin for to enjoy, then the Latin for to shape, then say \emph{queror, querī} again
\end{itemize}

\begin{tcolorbox}[breakable,colback=teal!4,colframe=teal!35!black,title={Before you move on}]
What does queror, querī mean? (\textquotedblleft{}To complain\textquotedblright{}.) Say it once more.
\end{tcolorbox}
\section[dēsum, deesse]{dēsum, deesse — to be missing, to fail}

\label{lesson:LA-C231-desum}



\begin{quote}
Before the new one: say the Latin for to shape, then the Latin for to complain.
\end{quote}

\subsection*{You'll want to know: dēsum, deesse}
\textbf{dēsum, deesse} — \textquotedblleft{}to be missing, to fail\textquotedblright{}.

\begin{cousinweb}
\textbf{dēsum} is \emph{dē}, \textquotedblleft{}away from\textquotedblright{}, with \emph{sum}, \textquotedblleft{}to be\textquotedblright{}: to be away, to be missing. \emph{Sum} is inherited from Proto-Indo-European \emph{*h\textsubscript{1}es-}, \textquotedblleft{}to be\textquotedblright{} (EtymDB 2.1).
\end{cousinweb}

\begin{grammarlens}[title={What you've built}]
One more word for this chapter.
\end{grammarlens}

\subsection*{Your turn}
\begin{itemize}
\raggedright
  \item \emph{Say it:} \emph{dēsum, deesse}
  \item \emph{Say it:} \emph{dēsum, deesse}, once more
  \item \emph{Say it:} say \emph{dēsum, deesse}
  \item \emph{Recall:} say the Latin for to shape, then the Latin for to complain, then say \emph{dēsum, deesse} again
\end{itemize}

\begin{tcolorbox}[breakable,colback=teal!4,colframe=teal!35!black,title={Before you move on}]
What does dēsum, deesse mean? (\textquotedblleft{}To be missing, to fail\textquotedblright{}.) Say it once more.
\end{tcolorbox}
